% Source: https://artofproblemsolving.com/community/c3h2752341_how_do_you_solve_these_factorization_problems
% Stable identity: aops-topic-2752341-q02
For example $a=1$, $b=2$, $c=3$ gives $\dfrac23+\dfrac45+\dfrac32=\dfrac{89}{30}\ne3$.

Since $\dfrac{2a}{a+b}=1+\dfrac{a-b}{a+b}$, the sum equals $3+T$ with $T=\dfrac{a-b}{a+b}+\dfrac{b-c}{b+c}+\dfrac{c-a}{c+a}$. Put everything over the common denominator $(a+b)(b+c)(c+a)$; expanding the numerator gives
\[T=-\dfrac{(a-b)(b-c)(c-a)}{(a+b)(b+c)(c+a)}.\]
(The numerator is a polynomial that vanishes when $a=b$, $b=c$ or $c=a$ and has degree $3$, so it is a constant multiple of $(a-b)(b-c)(c-a)$; comparing the coefficient of $a^2b$ gives the constant $-1$.) Hence the sum equals $3$ exactly when $(a-b)(b-c)(c-a)=0$, that is, when at least two of $a,b,c$ are equal. $\blacksquare$
